\section{CFU Effectiveness and Latency Trade-Offs}
\label{sec:disc_cfu}

The analysis in Section~\ref{sec:disc_bottleneck} indicated that inference at MNIST scale is largely memory-bound. This section looks at the clause scoring measurements in that context, and discusses the effect of the two custom instruction designs on inference performance.

The scoring comparison in Section~\ref{sec:results_scoring} isolates clause evaluation by measuring only the on-chip scoring phase of the per-class strategy. The three implementations differ by less than 2\%: CFU Design 1 is 0.85\% faster than the baseline (Section \ref{sec:inference_pipeline}), while CFU Design 2 is 0.91\% slower. Although the magnitudes are negligible, the differences are perfectly consistent across the test set, with Design 1 faster than the baseline on every one of the 200 samples and Design 2 slower on every sample. The differences are therefore real and reproducible rather than measurement noise, even though they are too small to matter in practice.

The fact that CFU Design 1 offers little improvement is understandable when looking at the work it offloads to hardware.  The baseline software already evaluates clause chunks using bitwise operations over 32-bit words, meaning a single chunk evaluation only needs the mask operations, comparisons and a branch. These are packed into a single instruction by CFU Design 1, however it seems the fixed minimum latency imposed on CFU instructions overshadows whatever time is saved, resulting in a very minimal time save.

The creation of CFU Design 2 was motivated by moving a larger portion of the inference pipeline into hardware, expecting an improvement of the already promising results found in the preceeding project thesis \cite{JB2025}. Testing however proved these expectations wrong, as the design was the slowest of the three. The reason for this is similar to the lackluster results of CFU Design 1: the unfavourable balance of the computation moved to hardware and the fixed minimum overhead of the custom instructions, in addition to the extra one cycle of latency added due to the second design being sequential. The second design adds two new custom instructions in addition to the one used in CFU Design 1, which commits a clause and retrieves class scores. These instructions move trivial integer operations to hardware, which are both few in number and inexpensive to execute in software. This explains why CFU Design 2 performs worse than the other implementations.

Even by settings aside which design is fastest, the benefit available from accelerating clause evaluation at this scale is inherently limited. In the per-class strategy, clause scoring account for only 80.32 ms of the 861.05 ms inference time, meaning even if this was reduced to zero it would only reduce the total by around 9\%. In the per-chunk strategy it is hard to decompose DDR2 reads from the total inference time, as mentioned in Section \ref{sec:results_experimental}, however we can assume there is a similar relationship between times taken by DDR2 reads and clause evaluation. In other words, no CFU design, however efficient, can address the bottleneck identified in Section \ref{sec:disc_bottleneck}, as the bottleneck lies in moving the model rather than evaluating it.

This stands in sharp contrast to the 4.42$\times$ speedup reported for the same approach on NoisyXOR in the preceding project thesis \cite{JB2025}. There are two main reasons to this. The first, and more fundamental, is the one explained above: at NoisyXOR scale the model was small enough to reside entirely in on-chip memory, so clause evaluation comprised a lot more of the time spent during inference, meaning an instruction evaluating clauses had a much bigger impact on performance. At the scale of this work, with the MNIST dataset and more clauses per class, the workload is instead dominated by movement of the model between on-chip and off-chip memory. The second, and likely even more important factor, is that the baseline implementation in this work is considerably more optimised than the one used previously. The baseline employed in this work implements some of the techniques used in \cite{zeng2025fastcompacttsetlinmachine}, using bitwise operations to evaluate 32 literals in a single operation, and early exit for early clause termination. In the baseline version in the preceding project thesis only evaluated a single literal at a time, which is considerably slower. In other words, part of the earlier speedup can be attributed to a weaker point of comparison.

It is worth emphasising that the lack of benefit is likely not a result of an expensive or poorly designed instruction. The custom logic added is small, with CFU Design~2, the more complex of the two designs, only adds 57 LUTs and 23 flip-flops over the baseline at the same 8 KB DMEM size (Table \ref{tab:res_processor}), which is a fraction of the total design. The instructions work correctly and are cheap to implement, however they target a part of the inference pipeline which is both fast and already optimized, meaning the possible performance gains are limited.




